%% bk-phys-faces.tex（bk-phys.sty が読む）
%% 物理シリーズ第2〜5巻の「4つの顔」のピクトグラム（第10回編集委員会、2026-10-18）
%% 第1巻と同じ決まり：たて・よこ −1〜1 の枠に描く。#1＝色、#2＝線の太さ。矢じりは ph arrow（大きさに合わせて縮む）。
%% ph@face@<巻>@<番号>＝大きな絵（表紙のタイル）、ph@facesm@<巻>@<番号>＝直径10 mm 以下の丸に入れる、形をへらした絵。
%% 部品は線と塗りだけ（白黒でも形で見分けられる）。中の字や記号は使わない（＋と−だけは形として描く）。

% 小さな絵の矢じり（塗った三角）  \ph@head{色}{先の点}{向き（度）}{長さ}  先から長さ L、幅 1.1L。線とは別に置く（指摘 No.4）
\def\ph@head#1#2#3#4{\fill[#1] #2 -- ++(#3+180-29:#4*1.14) -- ++(#3-90:#4*1.10) -- cycle;}

%% ============ 第2巻 熱学 ============
% ふくらむ：下から温めた気体がふくらみ、ピストンを押し上げる（たてのシリンダ・ピストン・中の気体・ほのお）
\def\ph@flame#1#2#3#4#5{% 色, x, 下のはし, 高さ, 幅（中に穴のあるほのお）
  \fill[#1, even odd rule]
    (#2,#3) .. controls (#2+0.62*#5,#3) and (#2+0.62*#5,#3+0.42*#4) .. (#2+0.30*#5,#3+0.66*#4)
      .. controls (#2+0.12*#5,#3+0.80*#4) and (#2+0.10*#5,#3+0.90*#4) .. (#2+0.04*#5,#3+#4)
      .. controls (#2-0.36*#5,#3+0.80*#4) and (#2-0.62*#5,#3+0.50*#4) .. (#2-0.56*#5,#3+0.28*#4)
      .. controls (#2-0.50*#5,#3+0.08*#4) and (#2-0.30*#5,#3) .. (#2,#3) -- cycle
    (#2,#3+0.10*#4) .. controls (#2+0.26*#5,#3+0.10*#4) and (#2+0.24*#5,#3+0.32*#4) .. (#2+0.06*#5,#3+0.50*#4)
      .. controls (#2-0.14*#5,#3+0.34*#4) and (#2-0.24*#5,#3+0.12*#4) .. (#2,#3+0.10*#4) -- cycle;}
\expandafter\def\csname ph@face@2@i\endcsname#1#2{%
  \draw[#1, line width=1.3*#2, line join=round] (-0.46,0.74) -- (-0.46,-0.30) -- (0.46,-0.30) -- (0.46,0.74);
  \fill[#1, rounded corners=0.03*\ph@u] (-0.38,0.20) rectangle (0.38,0.36);
  \draw[#1, line width=1.3*#2] (0,0.36) -- (0,0.66);
  \fill[#1] (-0.22,-0.14) circle (0.055) (0.16,-0.04) circle (0.055) (-0.04,0.06) circle (0.055) (0.24,-0.18) circle (0.055);
  \draw[#1, line width=#2, line cap=round, ph arrow] (0.76,-0.02) -- (0.76,0.66);
  \ph@flame{#1}{-0.20}{-0.96}{0.50}{0.36}\ph@flame{#1}{0.22}{-0.96}{0.50}{0.36}}
\expandafter\def\csname ph@facesm@2@i\endcsname#1#2{% 小さな丸では、ほのおを省き、ピストンが押し上げられる矢印だけにする
  \draw[#1, line width=#2, line join=miter] (-0.66,0.70) -- (-0.66,-0.80) -- (0.66,-0.80) -- (0.66,0.70);
  \fill[#1] (-0.56,-0.36) rectangle (0.56,-0.10);
  \draw[#1, line width=#2, line cap=butt, ph arrow sm] (0,-0.10) -- (0,0.98);}
% めぐる：p-V 線図の上を時計回りにめぐるサイクル（軸と、閉じた道すじ）
\expandafter\def\csname ph@face@2@ii\endcsname#1#2{%
  \draw[#1, line width=0.8*#2, line cap=round, line join=round] (-0.92,0.90) -- (-0.92,-0.84) -- (0.98,-0.84);
  \draw[#1, line width=1.4*#2, line join=round] (-0.56,0.66) .. controls (-0.14,0.36) and (0.30,0.16) .. (0.72,0.08)
    -- (0.72,-0.56) .. controls (0.30,-0.50) and (-0.14,-0.34) .. (-0.56,-0.08) -- cycle;
  \draw[#1, line width=1.4*#2, ph arrow] (0.72,0.08) -- (0.72,-0.30);
  \draw[#1, line width=1.4*#2, ph arrow] (-0.56,-0.08) -- (-0.56,0.42);}
\expandafter\def\csname ph@facesm@2@ii\endcsname#1#2{%
  \draw[#1, line width=0.7*#2, line cap=butt, line join=miter] (-0.96,0.96) -- (-0.96,-0.92) -- (0.98,-0.92);
  \draw[#1, line width=#2, line join=round] (-0.52,0.66) .. controls (-0.10,0.34) and (0.30,0.18) .. (0.70,0.10)
    -- (0.70,-0.58) .. controls (0.30,-0.52) and (-0.10,-0.36) .. (-0.52,-0.10) -- cycle;
  \ph@head{#1}{(0.70,-0.42)}{-90}{0.40}}
% 移る：熱い側から冷たい側へ、壁を通りぬけていく熱（ゆらぐ線）
\expandafter\def\csname ph@face@2@iii\endcsname#1#2{%
  \fill[#1] (-0.13,-0.86) rectangle (0.13,0.86);
  \foreach \y in {0.48,0,-0.48} {
    \draw[#1, line width=#2, line cap=round] (-0.96,\y) sin (-0.86,\y+0.08) cos (-0.76,\y) sin (-0.66,\y-0.08) cos (-0.56,\y)
      sin (-0.46,\y+0.08) cos (-0.36,\y) -- (-0.24,\y);
    \draw[#1, line width=#2, line cap=round, ph arrow] (0.24,\y) -- (0.34,\y) sin (0.44,\y+0.08) cos (0.54,\y) sin (0.64,\y-0.08) cos (0.74,\y) -- (0.98,\y);}}
\expandafter\def\csname ph@facesm@2@iii\endcsname#1#2{%
  \fill[#1] (-0.14,-0.90) rectangle (0.14,0.90);
  \draw[#1, line width=0.8*#2, line cap=butt] (-1.0,0) sin (-0.84,0.16) cos (-0.68,0) sin (-0.52,-0.16) cos (-0.36,0) -- (-0.26,0);
  \draw[#1, line width=0.8*#2, line cap=butt] (0.26,0) -- (0.56,0);
  \ph@head{#1}{(0.98,0)}{0}{0.46}}
% ぶつかる：箱の中を飛びまわる分子（点と、速さの矢印）。1つは右の壁にぶつかる
\expandafter\def\csname ph@face@2@iv\endcsname#1#2{% 1つは右の壁にぶつかってはね返る（来た道は点線）
  \draw[#1, line width=1.3*#2, rounded corners=0.08*\ph@u] (-0.88,-0.84) rectangle (0.88,0.84);
  \foreach \x/\y/\a in {-0.48/0.40/150, 0.02/0.52/40, -0.42/-0.36/250, 0.02/-0.10/190, 0.40/-0.56/250} {
    \fill[#1] (\x,\y) circle (0.10);
    \draw[#1, line width=0.8*#2, line cap=round, ph arrow] ($(\x,\y)+(\a:0.15)$) -- ($(\x,\y)+(\a:0.42)$);}
  \draw[#1, line width=0.6*#2, dash pattern=on 0.05*\ph@u off 0.05*\ph@u] (0.34,-0.16) -- (0.68,0.10);
  \fill[#1] (0.755,0.16) circle (0.10);
  \draw[#1, line width=0.8*#2, line cap=round, ph arrow] (0.66,0.26) -- (0.40,0.52);}
\expandafter\def\csname ph@facesm@2@iv\endcsname#1#2{% 小さな丸では、同じ線の上で向かいあう2つの分子にする
  \draw[#1, line width=0.8*#2] (-0.92,-0.88) rectangle (0.92,0.88);
  \fill[#1] (-0.60,0) circle (0.17) (0.60,0) circle (0.17);
  \draw[#1, line width=0.8*#2, line cap=butt] (-0.40,0) -- (-0.30,0) (0.40,0) -- (0.30,0);
  \ph@head{#1}{(-0.04,0)}{0}{0.30}\ph@head{#1}{(0.04,0)}{180}{0.30}}

%% ============ 第3巻 波動 ============
% ゆれる：天井からつるした、ばねとおもり。上下にゆれる（両向きの矢印）
\expandafter\def\csname ph@face@3@i\endcsname#1#2{%
  \draw[#1, line width=0.8*#2] (-0.68,0.86) -- (0.20,0.86);
  \foreach \x in {-0.62,-0.44,...,0.16} \draw[#1, line width=0.45*#2] (\x,0.86) -- (\x+0.08,0.98);
  \draw[#1, line width=0.8*#2, line join=round, line cap=round] (-0.24,0.86) -- (-0.24,0.76) -- (-0.42,0.69) -- (-0.06,0.57)
    -- (-0.42,0.45) -- (-0.06,0.33) -- (-0.42,0.21) -- (-0.24,0.14) -- (-0.24,0.06);
  \fill[#1, rounded corners=0.06*\ph@u] (-0.56,-0.56) rectangle (0.08,0.06);
  \draw[#1, line width=#2, line cap=round, ph arrow] (0.52,-0.18) -- (0.52,0.42);
  \draw[#1, line width=#2, line cap=round, ph arrow] (0.52,-0.18) -- (0.52,-0.78);}
\expandafter\def\csname ph@facesm@3@i\endcsname#1#2{%
  \draw[#1, line width=0.8*#2] (-0.80,0.94) -- (0.20,0.94);
  \draw[#1, line width=0.8*#2, line join=miter] (-0.30,0.94) -- (-0.54,0.76) -- (-0.06,0.56) -- (-0.54,0.36) -- (-0.30,0.20) -- (-0.30,0.10);
  \fill[#1] (-0.70,-0.62) rectangle (0.10,0.10);
  \draw[#1, line width=#2, line cap=butt, ph arrow sm] (0.62,-0.26) -- (0.62,0.58);
  \draw[#1, line width=#2, line cap=butt, ph arrow sm] (0.62,-0.20) -- (0.62,-1.0);}
% 伝わる：右へ進む波（正弦波と、進む向きの矢印）
\expandafter\def\csname ph@face@3@ii\endcsname#1#2{%
  \draw[#1, line width=0.5*#2, dash pattern=on 0.05*\ph@u off 0.05*\ph@u] (-0.96,-0.04) -- (0.96,-0.04);
  \draw[#1, line width=1.4*#2, line cap=round] (-0.96,-0.04) sin (-0.72,0.44) cos (-0.48,-0.04) sin (-0.24,-0.52) cos (0,-0.04)
    sin (0.24,0.44) cos (0.48,-0.04);
  \draw[#1, line width=#2, line cap=round, ph arrow] (0.10,0.72) -- (0.96,0.72);}
\expandafter\def\csname ph@facesm@3@ii\endcsname#1#2{%
  \draw[#1, line width=#2, line cap=butt] (-0.98,-0.20) sin (-0.68,0.30) cos (-0.38,-0.20) sin (-0.08,-0.70) cos (0.22,-0.20);
  \draw[#1, line width=#2, line cap=butt, ph arrow sm] (-0.30,0.66) -- (0.98,0.66);}
% はかる：ゆれの波形（上）を、周波数ごとの大きさ（下の棒）に分けて読む
\expandafter\def\csname ph@face@3@iii\endcsname#1#2{%
  \draw[#1, line width=0.9*#2, line cap=round, line join=round] plot[smooth, samples=61, domain=-0.92:0.92]
    (\x, {0.62+0.16*sin(\x*560)+0.09*sin(\x*1500+40)});
  \draw[#1, line width=0.8*#2, line cap=round] (-0.92,-0.78) -- (0.92,-0.78);
  \foreach \x/\h in {-0.66/0.40, -0.33/0.96, 0/0.30, 0.33/0.62, 0.66/0.20}
    \fill[#1] (\x-0.11,-0.78) rectangle (\x+0.11,-0.78+\h);}
\expandafter\def\csname ph@facesm@3@iii\endcsname#1#2{% 上に波形を残す（指摘 No.25）
  \draw[#1, line width=0.7*#2, line cap=butt] (-0.94,0.66) sin (-0.70,0.94) cos (-0.46,0.66) sin (-0.22,0.38) cos (0.02,0.66)
    sin (0.26,0.94) cos (0.50,0.66) sin (0.74,0.38) cos (0.94,0.62);
  \draw[#1, line width=0.8*#2] (-0.96,-0.86) -- (0.96,-0.86);
  \foreach \x/\h in {-0.62/0.52, 0/1.08, 0.62/0.72}
    \fill[#1] (\x-0.22,-0.86) rectangle (\x+0.22,-0.86+\h);}
% 曲がる：プリズムで光が曲がり、色に分かれる（底の側へ曲がる）
\expandafter\def\csname ph@face@3@iv\endcsname#1#2{%
  \fill[#1, opacity=0.25] (-0.10,0.70) -- (0.46,-0.42) -- (-0.66,-0.42) -- cycle;
  \draw[#1, line width=1.3*#2, line join=round] (-0.10,0.70) -- (0.46,-0.42) -- (-0.66,-0.42) -- cycle;
  \draw[#1, line width=#2, line cap=round] (-1.0,-0.04) -- (-0.38,0.14);
  \draw[#1, line width=0.8*#2, line cap=round] (-0.38,0.14) -- (0.20,0.10);
  \foreach \y in {-0.10,-0.36,-0.62}
    \draw[#1, line width=0.8*#2, line cap=round, ph arrow] (0.20,0.10) -- (0.98,\y);}
\expandafter\def\csname ph@facesm@3@iv\endcsname#1#2{%
  \draw[#1, line width=#2, line join=miter] (-0.18,0.76) -- (0.46,-0.52) -- (-0.82,-0.52) -- cycle;
  \draw[#1, line width=#2, line cap=butt] (-1.0,0.0) -- (-0.50,0.12);
  \draw[#1, line width=#2, line cap=butt, ph arrow sm] (0.15,0.12) -- (0.98,-0.40);}

%% ============ 第4巻 電磁気学 ============
% 引きあう：＋と−の電荷と、＋から−へ向かう電気力線
\expandafter\def\csname ph@face@4@i\endcsname#1#2{%
  \draw[#1, line width=1.2*#2] (-0.60,0) circle (0.32) (0.60,0) circle (0.32);
  \draw[#1, line width=1.2*#2, line cap=round] (-0.76,0) -- (-0.44,0) (-0.60,-0.16) -- (-0.60,0.16) (0.44,0) -- (0.76,0);
  \draw[#1, line width=0.9*#2, line cap=round, ph arrow] (-0.24,0) -- (0.24,0);
  \draw[#1, line width=0.9*#2, line cap=round, ph arrow] (-0.42,0.30) .. controls (-0.18,0.66) and (0.18,0.66) .. (0.42,0.30);
  \draw[#1, line width=0.9*#2, line cap=round, ph arrow] (-0.42,-0.30) .. controls (-0.18,-0.66) and (0.18,-0.66) .. (0.42,-0.30);}
\expandafter\def\csname ph@facesm@4@i\endcsname#1#2{%
  \draw[#1, line width=0.8*#2] (-0.58,-0.12) circle (0.38) (0.58,-0.12) circle (0.38);
  \draw[#1, line width=0.8*#2, line cap=butt] (-0.78,-0.12) -- (-0.38,-0.12) (-0.58,-0.32) -- (-0.58,0.08) (0.38,-0.12) -- (0.78,-0.12);
  \draw[#1, line width=0.8*#2, line cap=butt] (-0.50,0.30) .. controls (-0.28,0.86) and (0.20,0.90) .. (0.40,0.50);
  \ph@head{#1}{(0.52,0.28)}{-62}{0.40}}
% うず巻く：上へ流れる電流のまわりに、磁場がうずを巻く（上から見て反時計回り＝手前の側で右向き）
\expandafter\def\csname ph@face@4@ii\endcsname#1#2{%
  \foreach \y in {0.34,-0.36} {
    \draw[#1, line width=#2] (0,\y) ellipse (0.70 and 0.20);
    \draw[#1, line width=#2, ph arrow] (0.10,\y-0.20) -- (0.30,\y-0.193);}
  \draw[#1, line width=1.6*#2, line cap=round, ph arrow] (0,-0.92) -- (0,0.96);}
\expandafter\def\csname ph@facesm@4@ii\endcsname#1#2{%
  \draw[#1, line width=0.8*#2] (0,-0.05) ellipse (0.86 and 0.32);
  \ph@head{#1}{(0.66,-0.27)}{12}{0.40}
  \draw[#1, line width=#2, line cap=butt] (0,-0.98) -- (0,0.50);
  \ph@head{#1}{(0,0.98)}{90}{0.50}}
% 変える：変圧器（鉄心と、巻き数のちがう2つのコイル。左は4回、右は2回巻き）
\expandafter\def\csname ph@face@4@iii\endcsname#1#2{%
  \draw[#1, line width=2.6*#2] (-0.50,-0.66) rectangle (0.50,0.66);
  \foreach \y in {-0.36,-0.12,0.12,0.36}
    \draw[#1, line width=0.9*#2] (-0.50,\y) ellipse (0.26 and 0.085);
  \foreach \y in {-0.12,0.12}
    \draw[#1, line width=0.9*#2] (0.50,\y) ellipse (0.26 and 0.085);
  \draw[#1, line width=0.9*#2, line cap=round] (-0.76,0.36) -- (-0.98,0.36) (-0.76,-0.36) -- (-0.98,-0.36)
    (0.76,0.12) -- (0.98,0.12) (0.76,-0.12) -- (0.98,-0.12);}
\expandafter\def\csname ph@facesm@4@iii\endcsname#1#2{%
  \draw[#1, line width=1.2*#2] (-0.48,-0.76) rectangle (0.48,0.76);
  \foreach \y in {-0.36,0,0.36} \draw[#1, line width=0.7*#2] (-0.48,\y) ellipse (0.40 and 0.13);
  \foreach \y in {-0.18,0.18} \draw[#1, line width=0.7*#2] (0.48,\y) ellipse (0.40 and 0.13);}
% 届く：アンテナの塔から、電波が広がって届く
\expandafter\def\csname ph@face@4@iv\endcsname#1#2{%
  \begin{scope}[shift={(0.16,0)}]
  \draw[#1, line width=1.2*#2, line join=round, line cap=round] (-0.92,-0.88) -- (-0.62,0.30) -- (-0.32,-0.88);
  \draw[#1, line width=0.8*#2, line cap=round] (-0.85,-0.60) -- (-0.39,-0.60) (-0.78,-0.30) -- (-0.46,-0.30) (-0.71,0.0) -- (-0.53,0.0)
    (-0.85,-0.60) -- (-0.46,-0.30) (-0.39,-0.60) -- (-0.78,-0.30) (-0.78,-0.30) -- (-0.53,0.0) (-0.46,-0.30) -- (-0.71,0.0);
  \fill[#1] (-0.62,0.40) circle (0.10);
  \foreach \r in {0.40,0.76,1.12}
    \draw[#1, line width=#2, line cap=round] ([shift={(-0.62,0.40)}]-36:\r) arc (-36:36:\r);
  \end{scope}}
\expandafter\def\csname ph@facesm@4@iv\endcsname#1#2{%
  \draw[#1, line width=#2, line join=miter] (-0.80,-0.96) -- (-0.50,0.16) -- (-0.20,-0.96);
  \fill[#1] (-0.50,0.28) circle (0.17);
  \foreach \r in {0.62,1.18}
    \draw[#1, line width=#2, line cap=butt] ([shift={(-0.50,0.28)}]-38:\r) arc (-38:38:\r);}

%% ============ 第5巻 原子 ============
% 飛び出す：光（ゆらぐ線）が金属に当たり、電子（点）が飛び出す
\expandafter\def\csname ph@face@5@i\endcsname#1#2{%
  \fill[#1] (-0.96,-0.70) rectangle (0.96,-0.50);
  \foreach \x in {-0.88,-0.64,...,0.90} \draw[#1, line width=0.45*#2] (\x,-0.70) -- (\x-0.08,-0.82);
  \draw[#1, line width=#2, line cap=round, decorate, decoration={snake, segment length=0.22*\ph@u, amplitude=0.07*\ph@u, post length=0.16*\ph@u}, ph arrow]
    (-0.90,0.82) -- (-0.16,-0.44);
  \fill[#1] (0.16,-0.30) circle (0.11);
  \draw[#1, line width=#2, line cap=round, ph arrow] (0.28,-0.14) -- (0.86,0.66);}
\expandafter\def\csname ph@facesm@5@i\endcsname#1#2{%
  \fill[#1] (-0.98,-0.86) rectangle (0.98,-0.56);
  \draw[#1, line width=0.8*#2, line cap=butt, decorate, decoration={snake, segment length=0.36*\ph@u, amplitude=0.12*\ph@u, post length=0.10*\ph@u}]
    (-0.90,0.84) -- (-0.30,-0.50);
  \fill[#1] (0.10,-0.30) circle (0.19);
  \draw[#1, line width=#2, line cap=butt, ph arrow sm] (0.26,-0.08) -- (0.90,0.86);}
% すりぬける：波が壁（エネルギーの山）をすりぬけて、小さくなって向こうへ進む（トンネル効果）
\expandafter\def\csname ph@face@5@ii\endcsname#1#2{%
  \fill[#1] (-0.12,-0.66) rectangle (0.14,0.66);
  \draw[#1, line width=1.3*#2, line cap=round] (-0.96,0) sin (-0.83,0.46) cos (-0.70,0) sin (-0.57,-0.46) cos (-0.44,0)
    sin (-0.31,0.46) cos (-0.20,0.10);
  \draw[#1, line width=#2, line cap=round, ph arrow] (0.22,0.04) sin (0.34,0.17) cos (0.46,0) sin (0.58,-0.17) cos (0.70,0) -- (0.98,0);}
\expandafter\def\csname ph@facesm@5@ii\endcsname#1#2{%
  \fill[#1] (-0.16,-0.78) rectangle (0.16,0.78);
  \draw[#1, line width=#2, line cap=butt] (-0.98,0) sin (-0.75,0.56) cos (-0.52,0) sin (-0.29,-0.56) cos (-0.22,-0.45);
  \draw[#1, line width=#2, line cap=butt, ph arrow sm] (0.24,0) -- (0.98,0);}
% 切りかえる：半導体のチップと、0と1が切りかわる信号（方形波）（指摘 No.19）
\expandafter\def\csname ph@face@5@iii\endcsname#1#2{%
  \begin{scope}[shift={(-0.36,0)}, scale=0.78]
  \fill[#1, even odd rule, rounded corners=0.05*\ph@u] (-0.52,-0.52) rectangle (0.52,0.52) (-0.27,-0.27) rectangle (0.27,0.27);
  \fill[#1] (-0.17,-0.17) rectangle (0.17,0.17);
  \foreach \t in {-0.36,-0.12,0.12,0.36} {
    \draw[#1, line width=#2, line cap=butt] (-0.52,\t) -- (-0.80,\t) (0.52,\t) -- (0.80,\t) (\t,-0.52) -- (\t,-0.80) (\t,0.52) -- (\t,0.80);}
  \end{scope}
  \draw[#1, line width=#2, line cap=round, line join=round] (0.36,-0.24) -- (0.46,-0.24) -- (0.46,0.24) -- (0.62,0.24) -- (0.62,-0.24)
    -- (0.78,-0.24) -- (0.78,0.24) -- (0.98,0.24);}
\expandafter\def\csname ph@facesm@5@iii\endcsname#1#2{%
  \fill[#1] (-0.86,-0.40) rectangle (-0.06,0.40);
  \foreach \t in {-0.20,0.20} {
    \draw[#1, line width=0.7*#2, line cap=butt] (-0.86,\t) -- (-1.0,\t) (-0.06,\t) -- (0.08,\t) (\t-0.46,-0.40) -- (\t-0.46,-0.62) (\t-0.46,0.40) -- (\t-0.46,0.62);}
  \draw[#1, line width=0.8*#2, line cap=butt, line join=miter] (0.18,-0.30) -- (0.30,-0.30) -- (0.30,0.30) -- (0.62,0.30) -- (0.62,-0.30) -- (0.98,-0.30);}
% 分かれる：大きな原子核が2つに分かれ、中性子（小さな点）が飛び出す（核分裂）
\expandafter\def\csname ph@face@5@iv\endcsname#1#2{%
  \fill[#1] (-0.52,0) circle (0.38);
  \fill[#1] (0.64,0.50) circle (0.24) (0.64,-0.50) circle (0.24);
  \draw[#1, line width=#2, line cap=round, ph arrow] (-0.08,0.18) -- (0.34,0.40);
  \draw[#1, line width=#2, line cap=round, ph arrow] (-0.08,-0.18) -- (0.34,-0.40);
  \fill[#1] (0.16,0) circle (0.075);
  \draw[#1, line width=0.8*#2, line cap=round, ph arrow] (0.30,0) -- (0.96,0);}
\expandafter\def\csname ph@facesm@5@iv\endcsname#1#2{%
  \fill[#1] (-0.56,0) circle (0.42);
  \fill[#1] (0.66,0.58) circle (0.27) (0.66,-0.58) circle (0.27);
  \draw[#1, line width=0.8*#2, line cap=butt] (-0.10,0.06) -- (0.16,0.22) (-0.10,-0.06) -- (0.16,-0.22);
  \ph@head{#1}{(0.40,0.38)}{33}{0.36}\ph@head{#1}{(0.40,-0.38)}{-33}{0.36}}

\endinput
